%%%%%%%%%%%%%%%%%%%%%%%%%%%%%%%%%%%%%%%%%%%%%%%%%%%%%%%%%%%%%%%%%%%%%%%%%%
%%% Customizações do abnTeX2 para o Campus Crateús do IFCE    			%%%
%%% Instituto Federal de Educação, Ciência e Tecnologia do Ceará.   	%%%
%%%																	%%%
%%% Template disponível em: https://github.com/clodomirneto/IFCETeX2	%%%
%%% Desenvolvedor do IFCETeX2: Professor Clodomir Silva Lima Neto  	%%%
%%% Email para contato: clodomir.neto@ifce.edu.br						%%%
%%% 																	%%%
%%% Agradecimento: Thiago Nascimento 									%%%
%%% https://github.com/thiagodnf/uecetex2                              %%%
%%%%%%%%%%%%%%%%%%%%%%%%%%%%%%%%%%%%%%%%%%%%%%%%%%%%%%%%%%%%%%%%%%%%%%%%%%

%%%%%%%%%%%%%%%%%%%%%%%%%%%%%%%%%%%%%%%%%%%%%%%%%%%%%%%%%%%%%%%%%%%%%%%%%%
%%% 						Sequência do GitHub							%%%
%%%																	%%%
%%% git status + git add . + git commit + git push					%%%
%%%%%%%%%%%%%%%%%%%%%%%%%%%%%%%%%%%%%%%%%%%%%%%%%%%%%%%%%%%%%%%%%%%%%%%%%%

\documentclass[        
    a4paper,          % Tamanho da folha A4
    12pt,             % Tamanho da fonte 12pt
    chapter=TITLE,    % Todos os capitulos devem ter caixa alta
    %section=TITLE,    % Todas as secoes devem ter caixa alta
    oneside,          % Usada para impressao em apenas uma face do papel
    english,          % Hifenizacoes em ingles
    spanish,          % Hifenizacoes em espanhol
    brazil            % Ultimo idioma eh o idioma padrao do documento
]{abntex2}

% Importações de pacotes
\usepackage[utf8]{inputenc}                         % Acentuação direta
\usepackage[T1]{fontenc}                            % Codificação da fonte em 8 bits
\usepackage{graphicx}                               % Inserir figuras
\usepackage{amsfonts,amssymb,amsmath}             % Fonte e símbolos matemáticos
\usepackage{booktabs}                               % Comandos para tabelas
\usepackage{verbatim}                               % Texto é interpretado como escrito no documento
\usepackage{multirow,array}                        % Múltiplas linhas e colunas em tabelas
\usepackage{indentfirst}                            % Endenta o primeiro parágrafo de cada seção.
\usepackage{listings}                               % Utilizar codigo fonte no documento
\usepackage{xcolor}
\usepackage{microtype}                              % Para melhorias de justificação?
\usepackage[portuguese,ruled,lined]{algorithm2e}    % Escrever algoritmos
\usepackage{algorithmic}                            % Criar Algoritmos  
%\usepackage{float}                                  % Utilizado para criação de floats
\usepackage{amsgen}
\usepackage{lipsum}                                 % Usar a simulação de texto Lorem Ipsum
%\usepackage{titlesec}                               % Permite alterar os títulos do documento
\usepackage{tocloft}                                % Permite alterar a formatação do Sumário
\usepackage{etoolbox}                               % Usado para alterar a fonte da Section no Sumário
\usepackage[nogroupskip,nonumberlist,acronym]{glossaries}                % Permite fazer o glossario
\usepackage{caption}                                % Altera o comportamento da tag caption
\usepackage[alf,abnt-emphasize=bf,bibjustif, recuo=0cm,abnt-etal-cite=3,abnt-etal-list=0,abnt-etal-text=it]{abntex2cite}  % Citações padrão ABNT
%\usepackage[bottom]{footmisc}                      % Mantém as notas de rodapé sempre na mesma posição
%\usepackage{times}                                 % Usa a fonte Times
\usepackage{mathptmx}                               % Usa a fonte Times New Roman							
%\usepackage{lmodern}                               % Usa a fonte Latin Modern
%\usepackage{subfig}                                % Posicionamento de figuras
%\usepackage{scalefnt}                              % Permite redimensionar tamanho da fonte
%\usepackage{color, colortbl}                       % Comandos de cores
%\usepackage{lscape}                                % Permite páginas em modo "paisagem"
%\usepackage{ae, aecompl}                           % Fontes de alta qualidade
%\usepackage{picinpar}                              % Dispor imagens em parágrafos
%\usepackage{latexsym}                              % Símbolos matemáticos
%\usepackage{upgreek}                               % Fonte letras gregas
\usepackage{appendix}                               % Gerar o apendice no final do documento
\usepackage{paracol}                                % Criar paragrafos sem identacao
\usepackage{ifcetex2}		                        	  % Biblioteca com as normas do IFCE para trabalhos academicos
\usepackage{pdfpages}                               % Incluir pdf no documento
\usepackage{textcomp}
% Organiza e gera a lista de abreviaturas, simbolos e glossario
\makeglossaries
% Gera o Indice do documento
\makeindex

%TEOREMAS

\newtheorem{proposition}{Proposição}[chapter]
\newtheorem{theorem}{Teorema}[chapter]
\newtheorem{lemma}{Lema}[chapter]
\newtheorem{definition}{Definição}[chapter]
\newtheorem{exemplo}{Exemplo}[chapter]
\newtheorem{corollary}{Corolário}[chapter]
\newtheorem{exercicio}{Exercício}[chapter]

%PROVAS

\newenvironment{prova}[1][Prova]{\noindent\textbf{#1.} }{\hfill\rule{0.5em}{0.5em}}
\newenvironment{exm}{\noindent{\textbf{Exemplo:}}}{$\diamond$}
\newenvironment{dem}[1][Demonstra\c c\~ao]{\noindent\textbf{#1.} }{\hfill\rule{2mm}{2mm}}
%arquivo de template para os PUDS

%definição das variaveis das seções
\newcommand{\disciplina}{\def \disciplina}
\newcommand{\imprimirdisciplina}{\disciplina}

\newcommand{\codigo}{\def \codigo}
\newcommand{\imprimircodigo}{\codigo}

\newcommand{\cargaHorariaTotal}{\def \cargaHorariaTotal}
\newcommand{\imprimircargaHorariaTotal}{\cargaHorariaTotal}

\newcommand{\cargaHorariaPratica}{\def \cargaHorariaPratica}
\newcommand{\imprimircargaHorariaPratica}{\cargaHorariaPratica}

\newcommand{\cargaHorariaTeorica}{\def \cargaHorariaTeorica}
\newcommand{\imprimircargaHorariaTeorica}{\cargaHorariaTeorica}

\newcommand{\creditos}{\def \creditos}
\newcommand{\imprimircreditos}{\creditos}

\newcommand{\codigoPrerequisitos}{\def \codigoPrerequisitos}
\newcommand{\imprimircodigoPrerequisitos}{\codigoPrerequisitos}

\newcommand{\semestre}{\def \semestre}
\newcommand{\imprimirsemestre}{\semestre}

\newcommand{\nivel}{\def \nivel}
\newcommand{\imprimirnivel}{\nivel}

\newcommand{\codigoEquivalencias}{\def \codigoEquivalencias}
\newcommand{\imprimircodigoEquivalencias}{\codigoEquivalencias}

\newcommand{\ementa}{\def \ementa}
\newcommand{\imprimirementa}{\ementa}

\newcommand{\objetivo}{\def \objetivo}
\newcommand{\imprimirobjetivo}{\objetivo}

\newcommand{\programa}{\def \programa}
\newcommand{\imprimirprograma}{\programa}

\newcommand{\metodologiaEnsino}{\def \metodologiaEnsino}
\newcommand{\imprimirmetodologiaEnsino}{\metodologiaEnsino}

\newcommand{\recursos}{\def \recursos}
\newcommand{\imprimirrecursos}{\recursos}

\newcommand{\avaliacao}{\def \avaliacao}
\newcommand{\imprimiravaliacao}{\avaliacao}

\newcommand{\bibliografiaBasica}{\def \bibliografiaBasica}
\newcommand{\imprimirbibliografiaBasica}{\bibliografiaBasica}

\newcommand{\bibliografiaComplementar}{\def \bibliografiaComplementar}
\newcommand{\imprimirbibliografiaComplementar}{\bibliografiaComplementar}

\newcommand{\versao}{\def \versao}
\newcommand{\imprimirversao}{\versao}


%comando de impressão da estrutura
\newcommand{\imprimirPUD}{
%Cabeçalho do PUD
\begin{Spacing}{1}

\noindent \begin{minipage}{2.5cm}%
\includegraphics[scale=0.12]{logo-ifce}
\end{minipage}
\hspace{0.3cm}
\begin{minipage}{13cm}%
\centering INSTITUTO FEDERAL DE EDUCAÇÃO, CIÊNCIA E TECNOLOGIA DO CEARÁ- IFCE\\
CAMPUS JUAZEIRO DO NORTE\\
CURSO SUPERIOR EM AUTOMAÇÃO INDUSTRIAL\\
PROGRAMA DE UNIDADE DIDÁTICA – PUD\\
\end{minipage}%
\end{Spacing}

\begin{longtable}{|p{14cm}|}
%primeiro cabeçalho
\hline
\rowcolor{lightgray}
\multicolumn{1}{p{14cm}}{\textbf{Disciplina: \imprimirdisciplina}}\\
\hline
\endfirsthead

%cabeçalho
\hline
continuação PUD \imprimirdisciplina\\
\hline
\endhead

\hline
continua...\\
\hline
\endfoot

\hline
\rowcolor{lightgray}

\begin{tabular}{p{5.5 cm}| l}
coordenação & departamento pedagogico\\[16 ex]
\end{tabular}\\

\hline

\endlastfoot

%elementos
\textbf{Código:} \imprimircodigo\\


\textbf{Carga Horária } Teórica: \imprimircargaHorariaTeorica, Prática \imprimircargaHorariaPratica, Total: \imprimircargaHorariaTotal\\


\textbf{Número de créditos:} \imprimircreditos\\


\textbf{Código pré-requisitos:} \imprimircodigoPrerequisitos\\


\textbf{Semestre:} \imprimirsemestre\\


\textbf{Nível:} \imprimirnivel\\
\hline

\rowcolor{lightgray}
\multicolumn{1}{|p{14cm}|}{\textbf{Ementa}}\\
\hline
\multicolumn{1}{|p{14cm}|}{\imprimirementa}\\
\hline

\rowcolor{lightgray}
\multicolumn{1}{|p{14cm}|}{\textbf{Objetivo}}\\
\hline
\imprimirobjetivo\\
\hline

\rowcolor{lightgray}
\multicolumn{1}{|p{14cm}|}{\textbf{Programa}}\\
\hline
\imprimirprograma\\
\hline

\rowcolor{lightgray}
\multicolumn{1}{|p{14cm}|}{\textbf{Metodologia de ensino}}\\
\hline
\imprimirmetodologiaEnsino\\
\hline

\rowcolor{lightgray}
\multicolumn{1}{|p{14cm}|}{\textbf{Recursos}}\\
\hline
\imprimirrecursos\\
\hline

\rowcolor{lightgray}
\multicolumn{1}{|p{14cm}|}{\textbf{Avaliação}}\\
\hline
\imprimiravaliacao\\
\hline

\rowcolor{lightgray}
\multicolumn{1}{|p{14cm}|}{\textbf{Bibliografia básica}}\\
\hline
\imprimirbibliografiaBasica\\
\hline

\rowcolor{lightgray}
\multicolumn{1}{|p{14cm}|}{\textbf{Bibliografia complementar}}\\
\hline
\imprimirbibliografiaComplementar\\
\hline
\end{longtable}
\pagebreak
}
\begin{document}

\disciplina{Eletrônica digital 1}
\codigo{AUT2403}
\cargaHorariaTotal{80}
\cargaHorariaPratica{20}
\cargaHorariaTeorica{60}
\creditos{4}
\codigoPrerequisitos{-}
\semestre{1º}
\nivel{superior}
\ementa{
Métodos de conversão de um sistema de numeração (decimal, binário, octal e
hexadecimal) e suas operações (soma, subtração, multiplicação). Representação de Números decimais usando o código BCD. Compreender o propósito dos Códigos alfanuméricos, como o código ASCII. Compreender as Operações e funções lógicas básicas (AND, OR e NOT) e suas funções derivadas. Avaliar o Potencial da álgebra de Booleana (teoremas, propriedades e postulados) e mapa de Karnaugh na simplificação de circuito lógicos complexos. Conhecer as Características básicas de CI’s digitais TTL e CMOS. Analisar o Funcionamento de circuitos lógicos combinacionais. Compreender os Circuitos somadores e subtratores. e Projetar Projeto de circuitos lógicos simples.
}

\objetivo{
• Realizar conversões numéricas das bases decimal, octal, hexadecimal e binário para seuequivalente em qualquer outro sistema de numeração.\\
• Realizar as operações aritméticas nas bases decimal, hexadecimal, octal e binário.\\
• Desenhar e interpretar os símbolos de portas lógicas do padrão IEEE/ANSI.\\
• Implementar circuitos lógicos usando as portas básicas AND, OR e NOT.\\
• Executar os passos necessários para obter a forma mais simplificada de uma expressão lógica.\\
• Interpretar os estudos de casos na análise de defeitos de circuitos combinacionais.\\
• Usar somadores completos no projeto de somadores binários paralelos.\\
• Implementar circuitos lógicos combinacionais.\\
}

\programa{
• Códigos binários Sistemas de numeração\\
• Sistema ponderado, bases 10, 2, 8 e 16 Conversão entre bases Aritmética binária\\
• Adição binária Subtração binária Multiplicação binária Complemento de dez Complemento de dois\\
• Álgebra de Boole\\
• Variáveis lógicas Tabelas da verdade\\
• Funções de uma variável Funções de duas variáveis lógicas\\
• Funções lógicas básicas (OR, AND e NOT) Funções lógicas derivadas Portas lógicas\\
• Propriedades da Álgebra de Boole Teoremas de Morgan\\
• Diagramas de Venn\\
• Levantamento de expressões lógicas Síntese de Circuitos Lógicos\\
• Tabelas da verdade e soma de produtos\\
• Realização de expressões lógicas com portas AND, NAND, OR e NOT Análise de Circuitos Lógicos\\
• Circuitos integrados digitais\\
• Características da família CMOS Características da família TTL; Minimização de expressões lógicas: Mapas de Karnaugh\\
• Circuitos Somadores: Soma em complemento de 2; Soma em complemento de 1; Meio-Somadores Somadores Completos; Codificadores e Decodificadores: Conversores de códigos; Decodificador\\
BCD-7 segmentos;
}

\metodologiaEnsino{
Aulas expositivas.\\
Leitura e pesquisa.\\
Aulas práticas em laboratório de Informática - Simuladores.\\
Aulas práticas em laboratório – Sistemas Digitais.\\
Resolução de lista de exercícios.\\
Desenvolvimento de projetos: software e hardware.\\
}

\recursos{
Livros contidos na bibliografia.\\
Artigos.\\
Quadro e pincel.\\
Data-show.\\
Laboratório.\\
Computadores, dispositivos, equipamentos e softwares.\\
Lista de exercícios.\\
}

\avaliacao{
Avaliação escrita.\\
Práticas individuais e em grupo no laboratório.\\
Relatório de prática.\\
Avaliação de exercícios resolvidos.\\
Poderão ser inseridas outras avaliações durante o semestre.\\
}

\bibliografiaBasica{
• LOURENÇO, A. C.; CRUZ, E. C. A.; FERREIRA, S. R. e CHOURI, S. Jr. Circuitos
digitais: estude e use. São Paulo: Érica: 2007.\\
• TEIXEIRA, Hugo Tanzarella; TAVARES, Marley Fagundes; PEREIRA, Rodrigo Vinícius
Mendonça. Sistemas digitais. Londrina : Editora e Distribuidora Educacional S.A. 2017.\\
• WIDMER, Neal S.; MOSS, Gregory L.; TOCCI, Ronald J. Sistemas Digitais: princípios e aplicações. 12a. Edição. São Paulo: Pearson Education do Brasil, 2018.\\
}

\bibliografiaComplementar{
• IDOETA, Ivan Valeije; CAPUANO, Francisco Gabriel. Elementos de eletrônica digital. 39. ed. rev.atual. São Paulo: Érica, 2007.\\
• LEACH, Donald P. Eletrônica digital no laboratório. São Paulo: Makron Books,1993.\\
• MALVINO, A. P.; LEACH, D. P. Eletrônica digital: princípios e aplicações. McGraw- Hill, 1988.\\
• MENDONÇA, A.; ZELENOVSKY, R. Eletronica digital: curso prático e exercicios. 2a Edição. Rio de Janeiro: MZ EDITORA. 2007.\\
}

\imprimirPUD

\end{document}
